\documentclass[twocolumn, 10pt]{article}

\usepackage[utf8]{inputenc}
\usepackage[spanish]{babel}
\usepackage{amsmath, amssymb}
\usepackage{graphicx}
\usepackage{geometry}
\usepackage{hyperref}
\usepackage{booktabs}
\usepackage{caption}
\usepackage{subcaption}
\usepackage{float}
\usepackage{parskip}
\usepackage{cite}

\geometry{
  top=1.8cm, bottom=1.8cm,
  left=1.5cm, right=1.5cm,
  columnsep=0.6cm
}

\captionsetup{font=small, labelfont=bf}
\captionsetup[subfigure]{font=scriptsize, labelfont=bf}

\title{%
  \large\textbf{Análisis de Datos Económicos Geolocalizados}\\
  \textbf{mediante Persistencia Homológica:}\\
  \textbf{Servicios de Salud General en Iztapalapa, CDMX}
}

\author{%
  % Nombre Apellido$^{1}$, Nombre Apellido$^{2}$\\[4pt]
  % \small $^{1,2}$Tecnológico de Monterrey\\
  % \small \texttt{correo@tec.mx}
}

\date{Mayo 2026}

\begin{document}

\maketitle
\thispagestyle{empty}

% ═══════════════════════════════════════
\begin{abstract}
Este trabajo aplica persistencia homológica al análisis de los datos
del Directorio Estadístico Nacional de Unidades Económicas (DENUE 2025)
en la alcaldía Iztapalapa, Ciudad de México.
A partir de 563 establecimientos de salud general georreferenciados,
se construyen complejos simpliciales de Vietoris--Rips y Alpha
(equivalente homológico al complejo de \v{C}ech) para calcular la
persistencia en dimensiones $H_0$ (conectividad) y $H_1$ (vacíos de
cobertura).
Los resultados muestran que el $87.4\%$ son consultorios privados, con
distancia mediana entre vecinos de 161.71\,m.
Se identifican ciclos $H_1$ persistentes que señalan zonas con anillos
de cobertura pero sin servicio interior, constituyendo candidatos
prioritarios para la instalación de nuevas unidades públicas de salud.
\end{abstract}

% ═══════════════════════════════════════
\section{Introducción}

La distribución espacial de los servicios de salud determina el acceso
efectivo de la población a la atención médica.
Métodos clásicos de análisis geoespacial —radios de influencia,
índices de densidad— no capturan la \emph{forma global} de la red ni
sus vacíos estructurales a múltiples escalas.
La \textbf{persistencia homológica} permite extraer características
topológicas intrínsecas: componentes conectadas ($H_0$), ciclos ($H_1$)
y cavidades ($H_2$) que persisten ante variaciones de escala.

El presente estudio aplica este enfoque a los establecimientos de
salud general (consultorios, clínicas y hospitales, códigos SCIAN
621111--622112) en Iztapalapa, la alcaldía más poblada de la Ciudad
de México ($\approx$1.8 millones de habitantes).
El objetivo es identificar patrones de conectividad y vacíos
persistentes de cobertura que orienten decisiones de política pública
y estrategia de negocio en el sector salud.

% ═══════════════════════════════════════
\section{Metodología}

\subsection{Preparación de los datos}

Se utilizan los datos abiertos del DENUE 2025 (INEGI) \cite{denue},
filtrando los registros con \texttt{cve\_ent}$=9$ y
\texttt{cve\_mun}$=10$ (Iztapalapa).
Se seleccionaron seis códigos SCIAN de salud general.
Tras la limpieza, se obtienen $n=563$ puntos con coordenadas válidas.

La Figura~\ref{fig:descriptivo} muestra la composición del conjunto de
datos: el $79.6\%$ corresponde a consultorios privados ($n=448$),
seguido de clínicas públicas ($6.6\%$) y clínicas privadas ($5.0\%$).
El $87.4\%$ pertenece al sector privado.
La Figura~\ref{fig:mapa_puntos} muestra la distribución geográfica de
todos los establecimientos dentro del límite de la alcaldía.

\begin{figure}[H]
  \centering
  \includegraphics[width=\linewidth]%
    {outputs/figures/01_barras_tipo_establecimiento.jpg}
  \caption{Conteo de establecimientos por tipo. Dominio de consultorios
    generales privados ($n=448$, $79.6\%$).}
  \label{fig:descriptivo}
\end{figure}

\begin{figure}[H]
  \centering
  \includegraphics[width=\linewidth]%
    {outputs/figures/06_mapa_puntos_por_tipo.jpg}
  \caption{Distribución geográfica de los 563 establecimientos de salud
    general en Iztapalapa, coloreados por tipo.}
  \label{fig:mapa_puntos}
\end{figure}

Las coordenadas se proyectan al sistema UTM zona 14N (EPSG:32614) para
operar en metros.
Se calculan variables derivadas: sector (público/privado), distancia al
vecino más cercano (BallTree haversine) e indicador de aislamiento
(percentil 90, umbral $359\,\text{m}$).

\subsection{Construcción de complejos simpliciales}

Sea $X=\{x_1,\ldots,x_{563}\}\subset\mathbb{R}^2$ la nube proyectada
y centrada.
Se construyen dos familias de complejos:

\textbf{Vietoris--Rips} $\mathrm{VR}(r)$: $\sigma\subseteq X$ es un
símplice si $d(x_i,x_j)\leq r$ para todo par.
Se evalúan radios $r\in\{300,700,1200,2000\}$\,m para visualización
(submuestra $n=80$) y $r_{\max}=5{,}000$\,m para la persistencia
completa, usando \texttt{Ripser} \cite{ripser}.

La Figura~\ref{fig:vr_evolucion} ilustra la evolución del complejo al
crecer el radio: a $r=300$\,m sólo existen 15~aristas locales; a
$r=700$\,m emergen los primeros ciclos de corto alcance; a
$r=1{,}200$\,m se consolidan triángulos en las zonas más densas; a
$r=2{,}000$\,m la red alcanza 295~aristas y conecta la mayor parte
de la alcaldía.

\begin{figure}[H]
  \centering
  \begin{subfigure}[b]{0.48\linewidth}
    \includegraphics[width=\linewidth]%
      {outputs/figures/vr_radio_300m.jpg}
    \caption{$r=300$\,m — 15 aristas}
  \end{subfigure}
  \hfill
  \begin{subfigure}[b]{0.48\linewidth}
    \includegraphics[width=\linewidth]%
      {outputs/figures/vr_radio_700m.jpg}
    \caption{$r=700$\,m — ciclos emergentes}
  \end{subfigure}

  \vspace{4pt}

  \begin{subfigure}[b]{0.48\linewidth}
    \includegraphics[width=\linewidth]%
      {outputs/figures/vr_radio_1200m.jpg}
    \caption{$r=1{,}200$\,m — triángulos}
  \end{subfigure}
  \hfill
  \begin{subfigure}[b]{0.48\linewidth}
    \includegraphics[width=\linewidth]%
      {outputs/figures/vr_radio_2000m.jpg}
    \caption{$r=2{,}000$\,m — 295 aristas}
  \end{subfigure}
  \caption{Crecimiento del complejo Vietoris--Rips sobre submuestra
    $n=80$ en Iztapalapa. Los puntos rojos son establecimientos;
    las aristas en azul conectan pares con $d\leq r$; los triángulos
    azul claro indican 2-símplices.}
  \label{fig:vr_evolucion}
\end{figure}

\textbf{Complejo Alpha} $\mathrm{Alpha}(\alpha)$: subconjunto de la
triangulación de Delaunay que satisface $H_*(\mathrm{Alpha}(\alpha))
\cong H_*(\check{C}(r))$ \cite{edelsbrunner}, construido con
\texttt{GUDHI} \cite{gudhi}.
La relación de intercalación
$\mathrm{VR}(r)\subseteq\check{C}(r)\subseteq\mathrm{VR}(2r)$
garantiza equivalencia homológica.

La Figura~\ref{fig:complejos} compara ambos complejos a radio
equivalente: el VR a $r=1{,}000$\,m y el Alpha con
$\alpha^2\leq700^2$\,m$^2$; ambos comparten la misma homología
pero el Alpha requiere significativamente menos símplices.

\begin{figure}[H]
  \centering
  \begin{subfigure}[b]{0.48\linewidth}
    \includegraphics[width=\linewidth]%
      {outputs/figures/complejo_vietoris_rips_iztapalapa_1000m.jpg}
    \caption{Vietoris--Rips ($r=1000$\,m)}
  \end{subfigure}
  \hfill
  \begin{subfigure}[b]{0.48\linewidth}
    \includegraphics[width=\linewidth]%
      {outputs/figures/12_complejo_alpha_iztapalapa.jpg}
    \caption{Alpha / \v{C}ech ($\alpha^2\leq700^2$\,m$^2$)}
  \end{subfigure}
  \caption{Complejos simpliciales sobre submuestra $n=80$ con el límite
    de Iztapalapa. Aristas en azul, triángulos en azul claro,
    vértices en rojo.}
  \label{fig:complejos}
\end{figure}

\subsection{Persistencia homológica}

Se calcula la persistencia con coeficientes en $\mathbb{Z}/2\mathbb{Z}$
mediante el algoritmo de reducción de matrices \cite{edelsbrunner}.
Cada característica queda descrita por un par $(b,d)$
—nacimiento y muerte— visualizado con:

\begin{itemize}
  \item \textbf{Diagrama de barras} (\emph{barcode}): segmento $[b,d)$
        por característica, ordenado por persistencia $d-b$.
  \item \textbf{Diagrama de dispersión}: puntos $(b,d)$ donde la
        distancia a la diagonal indica la significancia.
\end{itemize}

La Figura~\ref{fig:barcodes} muestra los diagramas de barras para
$H_0$ y $H_1$, y la Figura~\ref{fig:scatter} el diagrama de dispersión
con los tres ciclos $H_1$ más persistentes anotados.

\begin{figure}[H]
  \centering
  \includegraphics[width=\linewidth]%
    {outputs/figures/13_barcodes_h0_h1.jpg}
  \caption{Diagramas de barras. Izquierda: $H_0$ (componentes
    conectadas). Derecha: $H_1$ (ciclos / vacíos de cobertura).
    Las barras más largas indican características más persistentes.}
  \label{fig:barcodes}
\end{figure}

\begin{figure}[H]
  \centering
  \includegraphics[width=\linewidth]%
    {outputs/figures/14_diagrama_persistencia_mejorado.jpg}
  \caption{Izquierda: diagrama de dispersión $H_0$ (azul) y $H_1$
    (naranja); C1--C3 son los ciclos más persistentes. Derecha:
    ranking de persistencias $H_1$.}
  \label{fig:scatter}
\end{figure}

% ═══════════════════════════════════════
\section{Resultados}

\subsection{Conectividad (\texorpdfstring{$H_0$}{H0})}

El barcode $H_0$ muestra 562 componentes finitas que se fusionan
progresivamente.
La mediana del radio de fusión es $161.71\,\text{m}$, indicando que
la mitad de los establecimientos tienen un vecino a distancia
caminable.
A $r\approx500\,\text{m}$ la red ya está mayoritariamente conectada;
a $r=1{,}000\,\text{m}$ quedan menos de 10 componentes.
Los 57 establecimientos aislados (umbral $359\,\text{m}$) se
concentran en la periferia sur-oriente.

\subsection{Vacíos de cobertura (\texorpdfstring{$H_1$}{H1})}
\label{sec:resultados-h1}

Los ciclos $H_1$ persistentes identifican \textbf{anillos de
establecimientos con vacíos interiores}: zonas donde la cobertura
circunda pero no penetra.
De los $110$ ciclos detectados, $14$ son significativos
($d-b>\mu+\sigma=317$\,m), con persistencia máxima de $1{,}377$\,m
(Figura~\ref{fig:vacios_h1}).
Los tres ciclos más persistentes (C1--C3, Figura~\ref{fig:scatter})
señalan brechas estructurales de acceso a salud, no atribuibles a
ruido topológico.

\begin{figure}[H]
  \centering
  \includegraphics[width=\linewidth]%
    {outputs/figures/17_vacios_cobertura_h1.jpg}
  \caption{Izquierda: nacimiento del ciclo vs.\ persistencia, con
    umbral de significancia $\mu+\sigma$. Derecha: histograma de
    persistencias $H_1$. Los $14$ ciclos a la derecha del umbral
    corresponden a vacíos topológicos significativos.}
  \label{fig:vacios_h1}
\end{figure}

\subsection{Comparación por sector}

\begin{table}[H]
\centering
\small
\caption{Estadísticas por sector.}
\label{tab:sector}
\begin{tabular}{lcc}
\toprule
\textbf{Métrica} & \textbf{Privado} & \textbf{Público} \\
\midrule
Establecimientos   & 492 (87.4\%) & 71 (12.6\%) \\
Dist. mediana (m)  & $\sim$155 & $\sim$800 \\
Aislados (p90)     & 47 & 10 \\
Persistencia $H_1$ & baja varianza & alta varianza \\
\bottomrule
\end{tabular}
\end{table}

La red pública, más dispersa, genera ciclos $H_1$ de mayor
persistencia individual —brechas de acceso más críticas— respecto a
la red privada, cuya mayor densidad produce ciclos más cortos y
numerosos.
La Figura~\ref{fig:aislados} localiza los $57$ establecimientos
aislados ($d>359$\,m al vecino más cercano), concentrados en la franja
sur-oriente de la alcaldía.

\begin{figure}[H]
  \centering
  \includegraphics[width=\linewidth]%
    {outputs/figures/10_mapa_establecimientos_aislados.jpg}
  \caption{Establecimientos aislados (estrellas) por tipo sobre la
    nube general (gris). La concentración periférica anticipa la
    localización de los huecos topológicos $H_1$.}
  \label{fig:aislados}
\end{figure}

% ═══════════════════════════════════════
\subsection{Propuesta de ubicación de nuevas unidades}
\label{sec:propuesta}

Para traducir el análisis topológico en una recomendación accionable
se construye una \emph{simulación de inserción}: dado que los ciclos
$H_1$ persistentes corresponden a anillos de cobertura con vacío
interior, el centro de cada hueco es un candidato natural para una
nueva unidad pública de salud.

\textbf{Procedimiento.}
(i)~Se construye una malla de $250\times250$ celdas sobre el
\emph{bounding box} de Iztapalapa y se conservan los $29{,}017$
puntos interiores al polígono.
(ii)~Para cada punto se calcula la distancia euclidiana al
establecimiento más cercano mediante \texttt{BallTree}.
(iii)~Se identifican los máximos locales de esta función (ventana de
$25$ celdas, umbral mínimo $400$\,m), que corresponden a los centros
geométricos de los huecos.
(iv)~Se seleccionan las $5$ ubicaciones con mayor déficit local
(Tabla~\ref{tab:propuesta}).
(v)~Se reinserta cada candidato en la nube y se recalcula la
persistencia con \texttt{Ripser} para evaluar su efecto.

\begin{table}[H]
\centering
\small
\caption{Cinco ubicaciones propuestas. ``Hueco'' = distancia al
  establecimiento más cercano (m).}
\label{tab:propuesta}
\begin{tabular}{clcc}
\toprule
\textbf{ID} & \textbf{Colonia (CP)} & \textbf{Hueco} & \textbf{Lat,\,Lon} \\
\midrule
P1 & Leyes de Reforma 3a~Sec.~(09310) & 1{,}561 & 19.390,\,$-$99.087 \\
P2 & Campestre Potrero (09637)        & 1{,}446 & 19.327,\,$-$98.991 \\
P3 & Lomas de la Estancia (09640)     & 1{,}191 & 19.325,\,$-$99.007 \\
P4 & Central de Abastos (09040)       & 1{,}169 & 19.377,\,$-$99.084 \\
P5 & San Juan Xalpa (09849)           &    989  & 19.343,\,$-$99.091 \\
\bottomrule
\end{tabular}
\end{table}

La Figura~\ref{fig:propuesta} sobrepone las cinco ubicaciones P1--P5
sobre un mapa de calor de la distancia al servicio más cercano.
Las zonas más rojas concentran los mayores déficits de cobertura
interna y coinciden con los centros de los ciclos $H_1$ persistentes
identificados en la Sección~\ref{sec:resultados-h1}.

\begin{figure}[H]
  \centering
  \includegraphics[width=\linewidth]%
    {outputs/figures/20_propuesta_ubicaciones.jpg}
  \caption{Propuesta de cinco nuevas unidades de salud (estrellas
    verdes P1--P5). Fondo: distancia al establecimiento más cercano;
    a mayor intensidad de rojo, mayor déficit local.}
  \label{fig:propuesta}
\end{figure}

\textbf{Efecto sobre la topología.}
La Figura~\ref{fig:simulacion} compara los diagramas de persistencia
antes ($n=563$) y después ($n=568$).
La persistencia máxima de $H_1$ se reduce de $1{,}377$\,m a $953$\,m
(\textbf{31\,\% menos}), confirmando que los cinco puntos propuestos
``cierran'' los ciclos más persistentes —es decir, eliminan los
vacíos estructurales de mayor escala.

\begin{figure}[H]
  \centering
  \includegraphics[width=\linewidth]%
    {outputs/figures/21_simulacion_antes_despues.jpg}
  \caption{Diagramas de persistencia antes y después de insertar las
    cinco ubicaciones propuestas. Los ciclos $H_1$ más altos se
    desplazan hacia la diagonal (menor persistencia): los huecos
    topológicamente más grandes desaparecen.}
  \label{fig:simulacion}
\end{figure}

\subsection{Validación territorial de vacíos topológicos}
\label{sec:validacion}

La persistencia homológica detecta huecos \emph{geométricos} en la
red, pero éstos no necesariamente representan brechas reales de
cobertura médica: un ciclo $H_1$ podría ubicarse sobre áreas verdes,
cerros, mercados mayoristas o grandes equipamientos sin demanda
residencial. Para evitar esta interpretación errónea se construye
una \textbf{máscara de actividad urbana} basada en el DENUE
\emph{completo} de CDMX (\textit{i.e.}, todos los giros económicos,
no sólo salud), con $n=88{,}187$ unidades económicas y $719$ giros
en Iztapalapa.

\textbf{Criterio.} Una celda del grid se considera
\emph{zona urbana activa} si su distancia al establecimiento DENUE
más cercano (cualquier giro) es $\leq r_{\text{hab}}$.
Se fija el umbral primario $r_{\text{hab}}=300$\,m (caminata de
$\sim$4\,min) y se incluye una verificación de sensibilidad a
$r_{\text{hab}}=500$\,m.
A $r_{\text{hab}}=300$\,m la máscara cubre el $98.0\,\%$ del
territorio interno de la alcaldía, confirmando que Iztapalapa es
mayoritariamente habitada; los huecos no cubiertos coinciden con
parques, zonas arqueológicas y la Sierra de Santa Catarina.

\textbf{Clasificación de los candidatos.} Para cada candidato
$P_i$ se aproxima su ``disco interior'' como un disco de radio
$r_{\text{disco}}=\tfrac{1}{2}\,(\text{persistencia del hueco})$
centrado en $P_i$. Se muestrean uniformemente $600$ puntos dentro
del disco y se calcula el porcentaje que cae en zona urbana activa.
La regla es:
\begin{itemize}\itemsep0pt
  \item $\geq 60\,\%$ urbano $\Rightarrow$ \textbf{accionable}
  \item $\leq 40\,\%$ urbano $\Rightarrow$ \textbf{falso positivo}
  \item $40\,\%$--$60\,\%$ $\Rightarrow$ \textbf{mixto / revisión}
\end{itemize}

\begin{table}[H]
\centering
\small
\caption{Validación territorial de los candidatos P1--P5. Las
  columnas \%U\texorpdfstring{\textsubscript{300}}{300} y \%U\texorpdfstring{\textsubscript{500}}{500} dan el
  porcentaje del interior dentro de zona urbana activa al umbral
  indicado.}
\label{tab:validacion}
\begin{tabular}{cccl}
\toprule
\textbf{ID} & \textbf{\%U\textsubscript{300}} & \textbf{\%U\textsubscript{500}} & \textbf{Clase (300\,m)} \\
\midrule
P1 & 100.0 & 100.0 & accionable \\
P2 & 37.3  & 76.6  & falso positivo \\
P3 & 43.4  & 88.4  & mixto \\
P4 & 100.0 & 100.0 & accionable \\
P5 & 50.5  & 99.7  & mixto \\
\bottomrule
\end{tabular}
\end{table}

La Figura~\ref{fig:validacion} muestra (izquierda) la máscara binaria
de actividad urbana con los cinco candidatos clasificados, y
(derecha) el gradiente continuo de distancia al DENUE más cercano
con la curva de nivel $r_{\text{hab}}=300$\,m. Los candidatos
sobre la franja sur-oriente (P2, P3) caen parcialmente sobre la
Sierra de Santa Catarina, lo que explica su clasificación como
falso positivo o mixto.

\begin{figure}[H]
  \centering
  \includegraphics[width=\linewidth]%
    {outputs/figures/22_validacion_territorial.jpg}
  \caption{Validación territorial. Izquierda: máscara binaria
    ($r_{\text{hab}}=300$\,m) sobre el DENUE completo; verde claro =
    zona urbana activa. Derecha: gradiente continuo de la distancia
    al DENUE más cercano. Estrella = accionable, cuadrado = mixto,
    \texttt{X} = falso positivo.}
  \label{fig:validacion}
\end{figure}

\textbf{Sensibilidad al umbral.}
La Figura~\ref{fig:sensibilidad} compara las máscaras a $300$ y
$500$\,m. Subir el umbral a $500$\,m cubre el $99.7\,\%$ del
territorio y reclasifica P2, P3, P5 como accionables; sólo el
interior periférico de la Sierra queda fuera. Conservadoramente,
adoptamos el umbral $300$\,m como criterio principal y reportamos
$500$\,m como cota superior.

\begin{figure}[H]
  \centering
  \includegraphics[width=\linewidth]%
    {outputs/figures/23_sensibilidad_rhab.jpg}
  \caption{Sensibilidad de la máscara urbana al umbral
    $r_{\text{hab}}$. A $300$\,m cubre $98\,\%$ del territorio
    interno; a $500$\,m cubre $99.7\,\%$.}
  \label{fig:sensibilidad}
\end{figure}

\subsection{Visualizaciones interactivas}

Como complemento al reporte estático, se generaron cuatro mapas
interactivos en formato HTML (Folium) que permiten explorar la red
con \emph{zoom}, capas conmutables y \emph{pop-ups} con metadatos:
\texttt{mapa\_interactivo\_tipos.html} (distribución por tipo con
\emph{clusters}),
\texttt{mapa\_heatmap\_densidad.html} (mapa de calor general y de
hospitales),
\texttt{mapa\_cobertura\_500m.html} (círculos de cobertura
$r=500$\,m por sector), y
\texttt{mapa\_aislados\_hospitales.html} (aislados $p_{90}$ y
hospitales georreferenciados).
Estos artefactos acompañan al reporte en la carpeta \texttt{outputs/}.

% ═══════════════════════════════════════
\section{Conclusiones y recomendaciones}

La persistencia homológica permite caracterizar estructuralmente la
red de servicios de salud de Iztapalapa más allá de métricas de
densidad. Los hallazgos cuantitativos principales son:

\begin{enumerate}
  \item \textbf{Conectividad.}
        La red alcanza conexión sustancial a $r\approx500$\,m, pero
        $57$ establecimientos ($10\,\%$) permanecen aislados más allá
        de $359$\,m (umbral $p_{90}$).

  \item \textbf{Vacíos estructurales.}
        Existen $14$ ciclos $H_1$ significativos
        ($d-b>\mu+\sigma=317$\,m); el mayor mide $1{,}377$\,m de
        diámetro y se ubica en el sur-oriente de la alcaldía.

  \item \textbf{Asimetría público--privado.}
        El sector privado ($87.4\,\%$ de las unidades) concentra la
        oferta en el centro-norte; el sector público es más disperso
        y deja los huecos $H_1$ más persistentes sin cubrir.
\end{enumerate}

\noindent\textbf{De ciclos persistentes a recomendaciones.}
Los ciclos $H_1$ persistentes deben interpretarse como candidatos
\emph{preliminares} de brechas de cobertura. Tras aplicar la
máscara de actividad urbana del DENUE completo
(Sección~\ref{sec:validacion}), únicamente los ciclos cuyo
interior coincide principalmente con zonas urbanas activas se
consideran \emph{huecos accionables}; los ubicados sobre la Sierra
de Santa Catarina, áreas verdes o equipamientos no residenciales
se clasifican como falsos positivos territoriales.

\noindent\textbf{Recomendación de política pública.}
Bajo el criterio conservador ($r_{\text{hab}}=300$\,m), los
candidatos \emph{accionables} son:
\begin{itemize}\itemsep0pt
  \item \textbf{P1 — Leyes de Reforma 3a Sec.} (CP 09310),
        hueco $1{,}561$\,m, interior $100\,\%$ urbano.
  \item \textbf{P4 — Central de Abastos} (CP 09040),
        hueco $1{,}169$\,m, interior $100\,\%$ urbano.
\end{itemize}
Los candidatos \emph{mixtos} P3 (Lomas de la Estancia) y P5
(San Juan Xalpa) requieren validación con datos de población AGEB
antes de comprometer inversión; al umbral $500$\,m ambos cruzan a
accionables. P2 (Campestre Potrero) se descarta como falso
positivo territorial por su solapamiento con la Sierra de Santa
Catarina.
La simulación de inserción con los cinco candidatos reduce la
persistencia máxima de $H_1$ en $31\,\%$ ($1{,}377\to953$\,m);
restringiendo a los dos accionables se obtiene una reducción
parcial, pero defendible territorialmente.

\noindent\textbf{Recomendación de negocio.}
Para un operador privado de consultorios o farmacias con servicio
médico, P1 y P4 son los mercados más sólidos: zonas con demanda
potencial dentro de un radio caminable ($<1$\,km), sin competencia
directa y plenamente urbanizadas. P4 destaca adicionalmente por la
conectividad vial generada por la Central de Abastos y el flujo
diario de trabajadores.

\noindent\textbf{Valor del enfoque topológico.}
A diferencia de los análisis de densidad poblacional —que sólo
identifican zonas \emph{de baja densidad}—, la persistencia $H_1$
detecta vacíos \emph{estructurales}: regiones rodeadas de oferta
pero sin acceso interno. La combinación de persistencia $H_1$ con
una máscara territorial DENUE permite focalizar la inversión
donde una sola unidad nueva tiene el mayor impacto marginal en la
cobertura de la red, distinguiendo huecos reales de huecos
geográficos sin demanda residencial.

% ═══════════════════════════════════════
\begin{thebibliography}{9}

\bibitem{ripser}
U.~Bauer, ``Ripser: efficient computation of Vietoris--Rips
persistence barcodes,'' \emph{J. Appl. Comput. Topol.}, vol.~5,
pp.~391--423, 2021.

\bibitem{edelsbrunner}
H.~Edelsbrunner y J.~Harer, \emph{Computational Topology: An
Introduction}. Amer. Math. Soc., 2010.

\bibitem{gudhi}
GUDHI Project, ``GUDHI User and Reference Manual,'' v3.12.0, 2024.
\url{https://gudhi.inria.fr/}

\bibitem{denue}
INEGI, ``Directorio Estadístico Nacional de Unidades Económicas
(DENUE) 2025.'' México, 2025.

\end{thebibliography}

\end{document}
